\documentclass[a4paper,12pt,twoside]{memoir}

% Castellano
\usepackage[spanish,es-tabla]{babel}
\selectlanguage{spanish}
\usepackage[utf8]{inputenc}
\usepackage[T1]{fontenc}
\usepackage{lmodern} % scalable font
\usepackage{microtype}
\usepackage{placeins}

\RequirePackage{booktabs}
\RequirePackage[table]{xcolor}
\RequirePackage{xtab}
\RequirePackage{multirow}

% Links
\PassOptionsToPackage{hyphens}{url}\usepackage[colorlinks]{hyperref}
\hypersetup{
	colorlinks,
	linkcolor={black!40!black},
	citecolor={blue!65!black},
	urlcolor={blue!80!black}
}

% Ecuaciones
\usepackage{amsmath}

%
\usepackage{fancybox}

% Rutas de fichero / paquete
\newcommand{\ruta}[1]{{\sffamily #1}}

% Párrafos
\nonzeroparskip

% Huérfanas y viudas
\widowpenalty100000
\clubpenalty100000

% Evitar solapes en el header
\nouppercaseheads

% Imagenes
\usepackage{graphicx}
\newcommand{\imagen}[2]{
	\begin{figure}[!h]
		\centering
		\includegraphics[width=0.9\textwidth]{#1}
		\caption{#2}\label{fig:#1}
	\end{figure}
	\FloatBarrier
}

\newcommand{\imagenTam}[3]{
	\begin{figure}[!h]
		\centering
		\includegraphics[width=#3\textwidth]{#1}
		\caption{#2}\label{fig:#1}
	\end{figure}
	\FloatBarrier
}

\newcommand{\imagenflotante}[2]{
	\begin{figure}%[!h]
		\centering
		\includegraphics[width=0.9\textwidth]{#1}
		\caption{#2}\label{fig:#1}
	\end{figure}
}



% El comando \figura nos permite insertar figuras comodamente, y utilizando
% siempre el mismo formato. Los parametros son:
% 1 -> Porcentaje del ancho de página que ocupará la figura (de 0 a 1)
% 2 --> Fichero de la imagen
% 3 --> Texto a pie de imagen
% 4 --> Etiqueta (label) para referencias
% 5 --> Opciones que queramos pasarle al \includegraphics
% 6 --> Opciones de posicionamiento a pasarle a \begin{figure}
\newcommand{\figuraConPosicion}[6]{%
  \setlength{\anchoFloat}{#1\textwidth}%
  \addtolength{\anchoFloat}{-4\fboxsep}%
  \setlength{\anchoFigura}{\anchoFloat}%
  \begin{figure}[#6]
    \begin{center}%
      \Ovalbox{%
        \begin{minipage}{\anchoFloat}%
          \begin{center}%
            \includegraphics[width=\anchoFigura,#5]{#2}%
            \caption{#3}%
            \label{#4}%
          \end{center}%
        \end{minipage}
      }%
    \end{center}%
  \end{figure}%
}

%
% Comando para incluir imágenes en formato apaisado (sin marco).
\newcommand{\figuraApaisadaSinMarco}[5]{%
  \begin{figure}%
    \begin{center}%
    \includegraphics[angle=90,height=#1\textheight,#5]{#2}%
    \caption{#3}%
    \label{#4}%
    \end{center}%
  \end{figure}%
}
% Para las tablas
\newcommand{\otoprule}{\midrule [\heavyrulewidth]}
%
% Nuevo comando para tablas pequeñas (menos de una página).
\newcommand{\tablaSmall}[5]{%
 \begin{table}
  \begin{center}
   \rowcolors {2}{gray!35}{}
   \begin{tabular}{#2}
    \toprule
    #4
    \otoprule
    #5
    \bottomrule
   \end{tabular}
   \caption{#1}
   \label{tabla:#3}
  \end{center}
 \end{table}
}

%
%Para el float H de tablaSmallSinColores
\usepackage{float}

%
% Nuevo comando para tablas pequeñas (menos de una página).
\newcommand{\tablaSmallSinColores}[5]{%
 \begin{table}[H]
  \begin{center}
   \begin{tabular}{#2}
    \toprule
    #4
    \otoprule
    #5
    \bottomrule
   \end{tabular}
   \caption{#1}
   \label{tabla:#3}
  \end{center}
 \end{table}
}

\newcommand{\tablaApaisadaSmall}[5]{%
\begin{landscape}
  \begin{table}
   \begin{center}
    \rowcolors {2}{gray!35}{}
    \begin{tabular}{#2}
     \toprule
     #4
     \otoprule
     #5
     \bottomrule
    \end{tabular}
    \caption{#1}
    \label{tabla:#3}
   \end{center}
  \end{table}
\end{landscape}
}

%
% Nuevo comando para tablas grandes con cabecera y filas alternas coloreadas en gris.
\newcommand{\tabla}[6]{%
  \begin{center}
    \tablefirsthead{
      \toprule
      #5
      \otoprule
    }
    \tablehead{
      \multicolumn{#3}{l}{\small\sl continúa desde la página anterior}\\
      \toprule
      #5
      \otoprule
    }
    \tabletail{
      \hline
      \multicolumn{#3}{r}{\small\sl continúa en la página siguiente}\\
    }
    \tablelasttail{
      \hline
    }
    \bottomcaption{#1}
    \rowcolors {2}{gray!35}{}
    \begin{xtabular}{#2}
      #6
      \bottomrule
    \end{xtabular}
    \label{tabla:#4}
  \end{center}
}

%
% Nuevo comando para tablas grandes con cabecera.
\newcommand{\tablaSinColores}[6]{%
  \begin{center}
    \tablefirsthead{
      \toprule
      #5
      \otoprule
    }
    \tablehead{
      \multicolumn{#3}{l}{\small\sl continúa desde la página anterior}\\
      \toprule
      #5
      \otoprule
    }
    \tabletail{
      \hline
      \multicolumn{#3}{r}{\small\sl continúa en la página siguiente}\\
    }
    \tablelasttail{
      \hline
    }
    \bottomcaption{#1}
    \begin{xtabular}{#2}
      #6
      \bottomrule
    \end{xtabular}
    \label{tabla:#4}
  \end{center}
}

%
% Nuevo comando para tablas grandes sin cabecera.
\newcommand{\tablaSinCabecera}[5]{%
  \begin{center}
    \tablefirsthead{
      \toprule
    }
    \tablehead{
      \multicolumn{#3}{l}{\small\sl continúa desde la página anterior}\\
      \hline
    }
    \tabletail{
      \hline
      \multicolumn{#3}{r}{\small\sl continúa en la página siguiente}\\
    }
    \tablelasttail{
      \hline
    }
    \bottomcaption{#1}
  \begin{xtabular}{#2}
    #5
   \bottomrule
  \end{xtabular}
  \label{tabla:#4}
  \end{center}
}



\definecolor{cgoLight}{HTML}{EEEEEE}
\definecolor{cgoExtralight}{HTML}{FFFFFF}

%
% Nuevo comando para tablas grandes sin cabecera.
\newcommand{\tablaSinCabeceraConBandas}[5]{%
  \begin{center}
    \tablefirsthead{
      \toprule
    }
    \tablehead{
      \multicolumn{#3}{l}{\small\sl continúa desde la página anterior}\\
      \hline
    }
    \tabletail{
      \hline
      \multicolumn{#3}{r}{\small\sl continúa en la página siguiente}\\
    }
    \tablelasttail{
      \hline
    }
    \bottomcaption{#1}
    \rowcolors[]{1}{cgoExtralight}{cgoLight}

  \begin{xtabular}{#2}
    #5
   \bottomrule
  \end{xtabular}
  \label{tabla:#4}
  \end{center}
}




\graphicspath{ {./img/} }

% Capítulos
\chapterstyle{bianchi}
\newcommand{\capitulo}[2]{
	\setcounter{chapter}{#1}
	\setcounter{section}{0}
	\setcounter{figure}{0}
	\setcounter{table}{0}
	\chapter*{#2}
	\addcontentsline{toc}{chapter}{#2}
	\markboth{#2}{#2}
}

% Apéndices
\renewcommand{\appendixname}{Apéndice}
\renewcommand*\cftappendixname{\appendixname}

\newcommand{\apendice}[1]{
	%\renewcommand{\thechapter}{A}
	\chapter{#1}
}

\renewcommand*\cftappendixname{\appendixname\ }

% Formato de portada
\makeatletter
\usepackage{xcolor}
\newcommand{\tutor}[1]{\def\@tutor{#1}}
\newcommand{\course}[1]{\def\@course{#1}}
\definecolor{cpardoBox}{HTML}{E6E6FF}

\definecolor{darkblue}{RGB}{11, 24, 115}

\def\maketitle{
  \null
  \thispagestyle{empty}
  % Cabecera ----------------
\noindent\includegraphics[width=\textwidth]{cabecera}\vspace{1cm}%
  \vfill
  % Título proyecto y escudo informática ----------------
  \colorbox{cpardoBox}{%
    \begin{minipage}{.8\textwidth}
      \vspace{.5cm}\Large
      \begin{center}
      \textbf{TFG del Grado en Ingeniería Informática}\vspace{.6cm}\\
      \textbf{\LARGE\@title{}}
      \end{center}
      \vspace{.2cm}
    \end{minipage}

  }%
  \hfill\begin{minipage}{.20\textwidth}
    \includegraphics[width=\textwidth]{escudoInfor}
  \end{minipage}
  \vfill
  % Datos de alumno, curso y tutores ------------------
  \begin{center}%
  {%
    \noindent\LARGE
    Presentado por \@author{}\\ 
    en Universidad de Burgos --- \@date{}\\
    Tutor: \@tutor{}\\
  }%
  \end{center}%
  \null
  \cleardoublepage
  }
\makeatother

\newcommand{\nombre}{Javier Martín Castro}
\newcommand{\nomtutor}{Pedro Renedo Fernández}

% Datos de portada
\title{RedELA \\Documentación Técnica}
\author{Javier Martín Castro}
\tutor{Pedro Renedo Fernández}
\date{\today}

\begin{document}

\maketitle



\cleardoublepage



%%%%%%%%%%%%%%%%%%%%%%%%%%%%%%%%%%%%%%%%%%%%%%%%%%%%%%%%%%%%%%%%%%%%%%%%%%%%%%%%%%%%%%%%



\frontmatter


\clearpage

% Indices
\tableofcontents

\clearpage

\listoffigures

\clearpage

\listoftables

\clearpage

\mainmatter

\appendix

\apendice{Plan de Proyecto Software}

\section{Introducción}
La finalidad inicial de este proyecto era ofrecer la oportunidad para mejorar la calidad de vida de las personas que han sido diagnosticadas con algún tipo de Esclerósis Múltiple (EM). Para ello se empezó a desarrollar una aplicación frontend que de manera amigable que daría soporte a aquellas personas que hayan decido utilizarla y padezcan la enfermedad.

Durante los primeros meses y durante el desarrollo, nos encontramos ante una serie de inconvenientes tanto técnicos como de comunicación con nuestro cliente que nos debía de servir con un cliente para obtener los datos que debíamos mostrar en la aplicación. Al trasmitir nuestras inquietudes al tutor {\nomtutor} tomamos la decisión de ponernos en contacto con la \textbf{Asociación de Castilla y León ELACyL} y tomamos la decisión de reconducir el desarrollo a personas que padecían esta patología.

En este apartado se valorarán y analizarán las la planificación temporal así como el estudio de viabilidad.

\section{Planificación temporal Redem}
Para la planificación temporal se tendrán en cuenta el tiempo que conlleva la curva de aprendizaje para poner en marcha el proyecto, así como las distintas funcionalidades que se van implementando según las necesidades aportadas por el cliente.
A través del tiempo se realizan sprints para poder desarrollar y añadir las funcionalidades correspondientes al proyecto.

Este proyecto estará alojado en Github \url{https://github.com/jmc1005/redem}

\begin{itemize}
\tightlist
\item Se aplicó una estrategia de desarrollo a partir de (\emph{sprints}) mediante iteraciones incrementales.
\item Se realizan \emph{sprints} que conllevan poco tiempo para la creación de funcionalidades.
\item Al finalizar el \emph{sprint} correspondiente se incorpora a la release que posteriormente formará parte del producto.
\item Se realizaban reuniones con el cliente para revisar el producto. En estas revisiones se aportan mejoras/modificaciones.
\item Tras la obtención de requisitos de cliente, se planifican las nuevas tareas y que deciden en que \emph{sprint} estarán asignadas.
\item Cada tarea incluida se estima en tiempo en un tablero \emph{scrum}.
\item Cada tarea estimada se priorizan dentro del tablero \emph{scrum}.
\end{itemize}

\subsection{Sprint 0 (19/10/2023)}
Se realiza reunión con el tutor {\nomtutor} para asignación del proyecto para realizar la parte frontend de una aplicación
para dar soporte a usuarios diagnoticados con Esclerósis Múltiples (EM).

En ella se comenta una idea inicial sobre que herramientas se van a emplear y la selección del framework para crear el proyecto.

\subsection{Sprint 1 (27/10/2023)}
Se realiza una primera toma de contacto con el cliente con los requisitos iniciales. En ella nos transmite que está muy ilusionado en 
que queramos aceptar el reto y poder dar a luz un aplicación que podrá facilitar la vida de mucha gente.

Para los objetivos iniciales nos informa que está interesado en crear una aplicación móvil para personas que hayan sido diagnosticadas
con EM y así poder aportar una solución en su día a día.

Le indicamos que tenemos intención en realizar la aplicación en un lenguaje que podría utilizarse en distintos dispositivos así como podría 
ser usada en distintos Sistemas Operativos. 

Inicialmente nos indica que una versión inicial estaría disponible para dispositivos con sistema operativo Android.

\subsection{Sprint 2 (28/10/2023 al 04/02/2024)}
Se realiza una primera toma de contacto con el framework Flutter después de analizar que ventajas e inconvenientes podríamos encontrarnos
para realizar la aplicación móvil solicitada por el cliente.

Tras tomar la decisión se empieza una formación inicial en el lenguaje de programación Dart para entender su funcionamiento y poder realizar
el desarrollo del proyecto.

\subsection{Sprint 3 (05/10/2023 al 11/02/2024)}
Se realiza la creación del proyecto con el framework Flutter y con el lenguaje de programación Dart. Además, se añaden las configuraciones iniciales
y se crea la página inicial.

\subsection{Sprint 4 (12/02/2024 al 18/02/2024)}
Se añaden las páginas de de Login y Registro al proyecto.

\subsection{Sprint 5 (19/02/2024 al 25/02/2024)}
Crear página detalle del usuario y home.

\subsection{Sprint 6 (26/02/2024 al 10/03/2024)}
Crear aplicación responsive para varios dispositivos (web, tablet, móvil)
Corregir bug de idioma
Enviar correo al cliente para resolver dudas a cerca de los requisitos del usuario.

\subsection{Sprint 7 (11/03/2024 al 17/03/2024)}
Actualizar datos usuario
Actualizar documentación
Crear página Administrador (tras las inquietudes transmitidas al tutor sobre la inconsistencia de datos, nos aconseja que nos creemos una página para que el administrador pueda realizar la gestión de síntomas, tipos de esclerosis múltiples,...)

\section{Planificación temporal RedELA}
Una vez tomada la decisión de enfocar el trabajo y el proyecto una vez nos pusimos en contacto con la Asociación de Esclerosis Lateral Amiotrófica de Castilla y León y nos informaron de los requisitos que necesitarían para crear una aplicación que facilitara la vida a pacientes con ELA nos planificamos de manera similar al proyecto inicial.

Este proyecto estará alojado en Github \url{https://github.com/jmc1005/red-ela}

\begin{itemize}
\tightlist
\item Se aplicó una estrategia de desarrollo a partir de (\emph{sprints}) mediante iteraciones incrementales.
\item Se realizan \emph{sprints} que conllevan poco tiempo para la creación de funcionalidades.
\item Al finalizar el \emph{sprint} correspondiente se incorpora a la release que posteriormente formará parte del producto.
\item Se realizaban reuniones con el cliente para revisar el producto. En estas revisiones se aportan mejoras/modificaciones.
\item Tras la obtención de requisitos de cliente, se planifican las nuevas tareas y que deciden en que \emph{sprint} estarán asignadas.
\item Cada tarea incluida se estima en tiempo en un tablero \emph{scrum}.
\item Cada tarea estimada se priorizan dentro del tablero \emph{scrum}.
\end{itemize}

\subsection{Sprint 0 (08/04/2024 - 14/04/2024)}
Analizar cambio rumbo del TFG.
Reunión con Asociación ELACyL para presentar idea del proyecto.
Analizar documento con requisitos iniciales enviado por la asociación
\begin{itemize}
\tightlist
\item Actores que interactúan (Gestor/a de Casos, Paciente, Cuidador principal)
\item Secciones de la aplicación
\item Fases entrega
\end{itemize}

\subsection{Sprint 1 (15/04/2024 - 21/04/2024)}
\begin{itemize}
\tightlist
\item Crear proyecto red-ela
\item Añadir configuración inicial
\item Añadir internacionalización, en esta primera versión sólo estará disponible en español.
\item Crear página registro/acceso
\item Crear página administrador
\end{itemize}

\subsection{Sprint 2 (22/04/2024 - 28/04/2024)}
\begin{itemize}
\tightlist 
\item Crear servicio autenticación usuario
\item Crear servicio gestión usuario
\item Crear página listado de usuario para la gestión de usuarios
\item Crear página detalle usuario
\end{itemize}

\subsection{Sprint 3 (29/04/2024 - 05/05/2024)}
\begin{itemize}
\tightlist 
\item Analizar flujo de los actores Paciente, Cuidador principal y Enfermera Gestora de Casos
\item Actualizar campos Detalle usuario
\end{itemize}

\subsection{Sprint 4 (06/05/2024 - 12/05/2024)}
\begin{itemize}
\tightlist 
\item Daily cliente.
\item Actualizar campos Detalle usuario
\end{itemize}

\subsection{Sprint 5 (13/05/2024 - 26/05/2024)}
\begin{itemize}
\tightlist 
\item Analizar crear OTP (One Time Password) para el registro de la aplicación.
\item Analizar enviar email con código de invitación para acceder página OTP (One Time Password)
\item Crear página OTP (One Time Password)
\item Crear funcionalidad OTP (One Time Password) con firebase a través del móvil
\item Crear servicio cifrar datos usuarios
\end{itemize}

\subsection{Sprint 6 (22/05/2024 - 02/06/2024)}
\begin{itemize}
\tightlist 
\item Separar en pantalla Admin los usuarios en Gestores de Casos, Pacientes y Cuidadores.
\item Enviar email al invitar a usuarios a regitrarse en la aplicación.
\item Actualizar Gestor de Casos
\item Actualizar Cuidador
\item Añadir pantalla Home
\item Añadir pantalla Citas
\end{itemize}

\subsection{Sprint 7 (03/06/2024 - 16/06/2024)}
\begin{itemize}
\tightlist 
\item Crear servicio gestión de citas
\item Añadir funcionalidad cancelar cita
\item Actualizar pantalla home
\item Ajustes envío email/invitación.
\item Añadir notificaciones push
\end{itemize}

\subsection{Sprint 8 (17/06/2024 - 30/06/2024)}
\begin{itemize}
\tightlist 
\item Corrección bugs
\item Envío notificaciones
\item Crear/Actualizar workflows para despliegue android/web
\item Crear cuenta desarrollador para despliegue Android
\end{itemize}

\section{Estudio de viabilidad}

\subsection{Viabilidad económica}
En este apartado vamos a realizar un supuesto sobre el coste que tendría la realización del proyecto en una empresa.

\subsubsection{Tiempo de desarrollo}
Esta aplicación al ser personalizada y con funcionalidades complejas vamos a indicar que el tiempo de desarrollo completo de la primera versión serán de unos 6 meses.

\subsubsection{Costes Desarrollador}
Suponemos que la aplicación ha sido realizada por un desarrollador con contrato temporal durante 6 meses y se considera que los gastos ocasionados serían:
\begin{table}[H]
\centering
\begin{tabular}{lr}
\toprule
Concepto	&	Coste   \\
\midrule
Salario bruto & 1.833,33 €\\
Complementos & 0,00 €\\
Contingencias Comunes (23,60\%) & 432,67 €\\
Accidentes de Trabajo y Enfermedades Profesionales (2,25\%) & 41,25 €\\
Desempleo (6,70\%) & 122,83 € \\
Formación Profesional (0,60\%) & 11,00 € \\
Fondo de garantía salarial (FOGASA) (0,20\%) & 3,67 €\\
\textbf{Coste del empleado} & \textbf{2.444,75 €}\\
\bottomrule
\end{tabular}
\caption{Costes realizar aplicación por un desarrollador.}
\label{costesdesarrollador}
\end{table}

\subsubsection{Costes Material}
Los costes para que el desarrollador pueda realizar la aplicación y además, se tiene en cuenta que el material que necesita ronda una amortización durante 5 años, son los siguientes:
\begin{table}[H]
\centering
\begin{tabular}{lcc}
\toprule
Concepto	&	Coste	&	Coste Amortizado  \\
\midrule
Portátil & 1.400 €	& 140 € \\
Cuenta desarrollador & 25 € & --\\
Dispositivo móvil android & 280 € & 28 €\\
Firebase & coste por uso & --\\
\textbf{Total} & \textbf{1705 €} & \textbf{168 €}\\
\bottomrule
\end{tabular}
\caption{Costes del material o subscripciones para realizar aplicación.}
\label{costesmaterial}
\end{table}

\subsubsection{Monetización}
Esta aplicación va a ser una aplicación gratuita para ayudar a pacientes que hayan sido diagnosticadas con la enfermedad ELA, de momento será gratuita. Sería la Asociación de Esclerosis Lateral Amiotrófica de Castilla y León la que tendría que considerar la manera de sacarle rentabilidad a la aplicación. Algunas de las opciones podrían ser:
\begin{itemize}
\tightlist 
\item Donaciones
\item Publicidad
\item Crowdfunding
\end{itemize}

\subsection{Viabilidad legal}
Para definir la viabilidad legal~\cite{viabilidadLegal} de nuestro proyecto tenemos primeramente que comprender que tan importante es para nuestro desarrollo entender que hay que tener en cuenta lo siguiente:
\begin{itemize}
\tightlist 
\item Debe ajustarse a las leyes y regulaciones.
\item Afecta al coste, calidad, reputación e incluso a la duración.
\item Durante el desarrollo del proyecto este tendrá la participación y colaboración de diversos colaboradores, en nuestro caso el equipo del proyecto, el cliente.
\end{itemize}

Nuestra aplicación, al estar relacionada con la sanidad pública debemos tener cuidado con la privacidad del paciente.

\subsubsection{Software}
Para Flutter~\cite{flutter} que es un framework para construir aplicaciones multiplataforma, la viabilidad legal se centra en:
\begin{itemize}
\tightlist 
\item La licencia Apache 2.0, para crear y redistribuir la aplicación de código abierto.
\item La protección del framework mediante la propiedad intelectual de Google.
\item Derechos de autor.
\end{itemize}

Para Firebase~\cite{firebase} que fue creado por Google contiene algunos puntos a tener en cuenta:
\begin{itemize}
\tightlist 
\item Está creado bajo la licencia Apache 2.0, permitiendo a los desarrolladores poder utilizar, modificar y redistribuir sin tener alguna restricción y de manera gratuita.
\item La propiedad intelectual pertenece a Google.
\item Al crear la cuenta para utilizar la plataforma se aceptan los términos del servicio.
\item Los desarrolladores cuentan con los derechos de autor sobre su aplicación.
\item Los desarrolladores están obligados a cumplir las políticas de privacidad marcadas por Google.
\end{itemize}


\apendice{Especificación de Requisitos}

\section{Introducción}

En este punto vamos a definir los requisitos que va a tener nuestra primera versión del proyecto. Para ello vamos a seguir la \textbf{especificación de requisitos de sofware (ERS)}~\cite{wiki:ieee830-1998}, donde se recogen una serie de buenas prácticas y son las siguientes:
\begin{itemize}
\tightlist 
\item \textbf{Completa}. Todos los requerimientos deben estar contemplados en ella y deben estar definidas todas las referencias .
\item \textbf{Consistente}. Debe tener coherencia con los requerimientos y además, con otras especificaciones.
\item \textbf{Inequívoca}. Claridad a la hora de redactar el documento para que no haya malas interpretaciones.
\item \textbf{Correcta}. Se deben cumplir con los requisitos recogidos en la especificación.
\item \textbf{Trazable}. Hace referencia a la posibilidad de comprobar la historia, aplicación o ubicación de un item mediante su identificación almacenada y documentada.
\item \textbf{Priorizable}. Posibilidad de organizar los items según la relevancia y separándolos en esenciales, condicionales y opcionales.
\item \textbf{Modificable}. Dar la posibilidad de poder ser modificable.
\item \textbf{Verificable}. Tener la posibilidad de realizar las pruebas de manera finita y sin costo.
\item \textbf{Clara}. Debe ser entendible.
\end{itemize}

\section{Objetivos generales}
Nuestro proyecto tendrá como objetivos principales:
\begin{itemize}
\tightlist 
\item Crear una aplicación multiplataforma en Android y Web.
\item Gestionar las citas a través de los gestores de casos y asignárselas a pacientes con ELA.
\item Visualizar las citas que tengan en el día.
\item Visualizar las citas que tengan en un calendario.
\item Gestionar aquella información necesaria para los pacientes, como tratamientos, hospitales.
\end{itemize}
\section{Catálogo de requisitos}
Vamos a hacer en el catálogo de requisitos una mención de los requisitos específicos tras el análisis de los objetivos generales.

\subsection{Requisitos Funcionales (RF)}
\subsubsection{RF-1 Gestión Gestor Casos}
\begin{itemize}
\tightlist 
\item \textbf{RF-1.1 Invitar gestor casos}. El gestor de casos puede enviar una invitación para que un gestor de casos se registre en la aplicación.
\item \textbf{RF-1.2 Invitar paciente}. El gestor de casos puede enviar una invitación para que el paciente se registre en la aplicación.
\item \textbf{RF-1.3 Eliminar paciente}. El paciente podrá ser eliminado por parte del gestor de casos.
\item \textbf{RF-1.4 Listar pacientes}. El gestor de casos podrá ver un listado de pacientes asignados a él.
\item \textbf{RF-1.5 Ver paciente}. El gestor de casos podrá ver los datos del paciente seleccionado.
\item \textbf{RF-1.6 Editar gestor casos}. Podrá actualizar sus datos.
\item \textbf{RF-1.7 Ver citas}. Ver citas.
\item \textbf{RF-1.8 Crear cita}. Podrá crear cita.
\item \textbf{RF-1.9 Editar cita}. Podrá crear cita.
\item \textbf{RF-1.10 Cancelar cita}. Podrá cancelar cita.
\end{itemize}

\subsubsection{RF-2 Gestión Paciente}
\begin{itemize}
\tightlist 
\item \textbf{RF-2.1 Editar paciente}. El paciente será el encargado de gestionar sus datos.
\item \textbf{RF-2.2 Invitar cuidador principal}. El paciente puede enviar una invitación para que el cuidador se registre en la aplicación.
\item \textbf{RF-2.3 Ver usuarios asociados}. El paciente podrá ver la información del cuidador y el gestor de casos asociados.
\item \textbf{RF-2.4 Cancelar cita}. Podrá cancelar cita.
\end{itemize}

\subsubsection{RF-3 Gestión Cuidador Principal}
\begin{itemize}
\tightlist 
\item \textbf{RF-3.1 Editar cuidador}. El paciente será el encargado de gestionar sus datos.
\item \textbf{RF-3.2 Invitar cuidador principal}. El paciente puede enviar una invitación para que el cuidador se registre en la aplicación.
\item \textbf{RF-3.3 Ver pacientes asociados}. El cuidador podrá ver la información del paciente.
\item \textbf{RF-3.4 Cancelar cita}. Podrá cancelar cita.
\end{itemize}

\subsubsection{RF-4 Gestión Citas}
\begin{itemize}
\tightlist 
\item \textbf{RF-4.1 Crear cita}. Crear una cita nueva.
\item \textbf{RF-4.2 Cancelar cita}. Eliminar cita.
\item \textbf{RF-4.3 Listar citas}. Ver las citas creadas.
\item \textbf{RF-4.3 Editar cita}. Editar la cita.
\end{itemize}

\subsubsection{RF-5 Gestión Admin}
\begin{itemize}
\tightlist 
\item \textbf{RF-5.1 Invitar admin}. Envía invitación para que un usuario pueda registrarse como admin.
\item \textbf{RF-5.2 Invitar gestor casos}. Envía invitación para que un usuario pueda registrarse como gestor de casos.
\item \textbf{RF-5.3 Listar admin}. Puede ver los administradores de la aplicación.
\item \textbf{RF-5.4 Listar gestores de casos}. Puede ver los gestores de casos de la aplicación.
\item \textbf{RF-5.5 Listar pacientes}. Puede ver los pacientes de la aplicación.
\item \textbf{RF-5.6 Listar cuidadores}. Puede ver los cuidadores de la aplicación.
\item \textbf{RF-5.7 Eliminar usuario}. Puede eliminar usuario.
\item \textbf{RF-5.8 Ver usuario}. Ver usuario.
\item \textbf{RF-5.9 Listar hospitales}. Ver las hospitales.
\item \textbf{RF-5.10 Crear hospital}. Crear hospital.
\item \textbf{RF-5.11 Editar hospital}. Editar hospital.
\item \textbf{RF-5.12 Eliminar hospital}. Eliminar hospital.
\item \textbf{RF-5.13 Listar tratamientos}. Ver las tratamientos.
\item \textbf{RF-5.14 Crear tratamiento}. Crear tratamiento.
\item \textbf{RF-5.16 Editar tratamiento}. Editar tratamiento.
\item \textbf{RF-5.17 Eliminar tratamiento}. Eliminar tratamiento.
\item \textbf{RF-5.18 Editar admin}. Editar admin.
\end{itemize}

\section{Especificación de requisitos}
\subsection{Diagramas de flujo}
\imagen{admin-flujo}{Diagrama flujo admin}
\imagen{gestor-casos-flujo}{Diagrama flujo gestor casos}
\imagen{paciente-cuidador-flujo}{Diagrama flujo paciente/cuidador}
\subsection{Casos de Uso}
\begin{table}[p]
	\centering
	\begin{tabularx}{\linewidth}{ p{0.21\columnwidth} p{0.71\columnwidth} }
		\toprule
		\textbf{CU-1}    & \textbf{Gestión usuarios}\\
		\toprule
		\textbf{Versión}              & 1.0    \\
		\textbf{Autor}                & {\nombre}\\
		\textbf{Requisitos asociados} & RF-1, RF-2, RF-3, RF-5\\
		\textbf{Descripción}          & Permite gestionar los gestores de casos \\
		\textbf{Precondición}         & Firebase está disponible \\
		\textbf{Acciones}             &
		\begin{enumerate}
			\def\labelenumi{\arabic{enumi}.}
			\tightlist
			\item Accede a la aplicación
			\item Ver información personal
			\item Todos menos el admin ver las citas
			\item Todos menos el admin ver usuarios asociados
			\item Si es gestor o admin ver botón invitar a usuario
			\item Solo admin ver usuarios (administradores, gestores casos, pacientes y cuidadores principales)
			\item Solo admin ver tratamientos
			\item Solo admin ver hospitales
		\end{enumerate}\\
		\textbf{Postcondición}        & Los pacientes y citas coinciden con las que hay en la base de datos \\
		\textbf{Excepciones}          & Error al obtener los pacientes o citas \\
		\textbf{Importancia}          & Alta \\
		\bottomrule
	\end{tabularx}
	\caption{CU-1 Gestión Gestores de casos.}
\end{table}

\begin{table}[p]
	\centering
	\begin{tabularx}{\linewidth}{ p{0.21\columnwidth} p{0.71\columnwidth} }
		\toprule
		\textbf{CU-2}    & \textbf{Enviar invitación}\\
		\toprule
		\textbf{Versión}              & 1.0    \\
		\textbf{Autor}                & {\nombre} \\
		\textbf{Requisitos asociados} & RF-1.1, RF-1.2, RF-2.2, RF-3.2, RF-5.1, RF-5.2 \\
		\textbf{Descripción}          & Permite invitar a un usuario \\
		\textbf{Precondición}         & Firebase está disponible \\
		\textbf{Acciones}             &
		\begin{enumerate}
			\def\labelenumi{\arabic{enumi}.}
			\tightlist
			\item Accede a la aplicación
			\item Ir a la pestaña de usuarios
			\item Pulsar el botón (+) de añadir
			\item Sale un pop up para rellenar los datos
			\item Enviar invitación
		\end{enumerate}\\
		\textbf{Postcondición}        & La invitación es recibida por el usuario final \\
		\textbf{Excepciones}          & Error al enviar invitación \\
		\textbf{Importancia}          & Alta \\
		\bottomrule
	\end{tabularx}
	\caption{CU-2 Enviar invitación.}
\end{table}

\begin{table}[p]
	\centering
	\begin{tabularx}{\linewidth}{ p{0.21\columnwidth} p{0.71\columnwidth} }
		\toprule
		\textbf{CU-3}    & \textbf{Eliminar usuario}\\
		\toprule
		\textbf{Versión}              & 1.0    \\
		\textbf{Autor}                & {\nombre} \\
		\textbf{Requisitos asociados} & RF-1.3, RF-5.7 \\
		\textbf{Descripción}          & Permite eliminar un usuario \\
		\textbf{Precondición}         & Firebase está disponible \\
		\textbf{Acciones}             &
		\begin{enumerate}
			\def\labelenumi{\arabic{enumi}.}
			\tightlist
			\item Accede a la aplicación
			\item Ir a la pestaña de usuarios
			\item Pulsar papelera para borrar
		\end{enumerate}\\
		\textbf{Postcondición}        & El usuario se borra \\
		\textbf{Excepciones}          & Error al borrar \\
		\textbf{Importancia}          & Alta \\
		\bottomrule
	\end{tabularx}
	\caption{CU-3 Eliminar usuario.}
\end{table}

\begin{table}[p]
	\centering
	\begin{tabularx}{\linewidth}{ p{0.21\columnwidth} p{0.71\columnwidth} }
		\toprule
		\textbf{CU-4}    & \textbf{Listar usuarios}\\
		\toprule
		\textbf{Versión}              & 1.0    \\
		\textbf{Autor}                & {\nombre} \\
		\textbf{Requisitos asociados} & RF-1.4, RF-5.3, RF-5.4, RF-5.5, RF-5.6 \\
		\textbf{Descripción}          & Permite listar usuarios \\
		\textbf{Precondición}         & Firebase está disponible \\
		\textbf{Acciones}             &
		\begin{enumerate}
			\def\labelenumi{\arabic{enumi}.}
			\tightlist
			\item Accede a la aplicación
			\item Ir a la pestaña de usuarios
		\end{enumerate}\\
		\textbf{Postcondición}        & Se muestran los usuarios \\
		\textbf{Excepciones}          & Error al obtener listado \\
		\textbf{Importancia}          & Alta \\
		\bottomrule
	\end{tabularx}
	\caption{CU-4 Listar usuarios.}
\end{table}

\begin{table}[p]
	\centering
	\begin{tabularx}{\linewidth}{ p{0.21\columnwidth} p{0.71\columnwidth} }
		\toprule
		\textbf{CU-5}    & \textbf{Ver usuario}\\
		\toprule
		\textbf{Versión}              & 1.0    \\
		\textbf{Autor}                & {\nombre} \\
		\textbf{Requisitos asociados} & RF-1.5, RF-2.3, RF-3.3, RF-5.8 \\
		\textbf{Descripción}          & Permite ver un usuario \\
		\textbf{Precondición}         & Firebase está disponible \\
		\textbf{Acciones}             &
		\begin{enumerate}
			\def\labelenumi{\arabic{enumi}.}
			\tightlist
			\item Accede a la aplicación
			\item Ir a la pestaña de usuarios
			\item Pulsar en usuario para ver la información
		\end{enumerate}\\
		\textbf{Postcondición}        & El usuario se muestra con la información correcta \\
		\textbf{Excepciones}          & Error al obtener los datos\\
		\textbf{Importancia}          & Alta \\
		\bottomrule
	\end{tabularx}
	\caption{CU-5 Ver usuario.}
\end{table}

\begin{table}[p]
	\centering
	\begin{tabularx}{\linewidth}{ p{0.21\columnwidth} p{0.71\columnwidth} }
		\toprule
		\textbf{CU-6}    & \textbf{Editar usuario}\\
		\toprule
		\textbf{Versión}              & 1.0    \\
		\textbf{Autor}                & {\nombre} \\
		\textbf{Requisitos asociados} & RF-1.6, RF-2.1, RF-3.1, RF-5.8, RF-5.18 \\
		\textbf{Descripción}          & Permite editar un usuario \\
		\textbf{Precondición}         & Firebase está disponible \\
		\textbf{Acciones}             &
		\begin{enumerate}
			\def\labelenumi{\arabic{enumi}.}
			\tightlist
			\item Accede a la aplicación
			\item Pulsar en el icono usuario
			\item Editar la información
		\end{enumerate}\\
		\textbf{Postcondición}        & El usuario se borra \\
		\textbf{Excepciones}          & Error al borrar \\
		\textbf{Importancia}          & Alta \\
		\bottomrule
	\end{tabularx}
	\caption{CU-6 Editar usuario.}
\end{table}

\begin{table}[p]
	\centering
	\begin{tabularx}{\linewidth}{ p{0.21\columnwidth} p{0.71\columnwidth} }
		\toprule
		\textbf{CU-7}    & \textbf{Gestión citas}\\
		\toprule
		\textbf{Versión}              & 1.0    \\
		\textbf{Autor}                & {\nombre} \\
		\textbf{Requisitos asociados} & RF-1.8, RF-1.9, RF-1.10, RF-2.4, RF-3.4, RF-4 \\
		\textbf{Descripción}          & Permite ver las citas\\
		\textbf{Precondición}         & Firebase está disponible \\
		\textbf{Acciones}             &
		\begin{enumerate}
			\def\labelenumi{\arabic{enumi}.}
			\tightlist
			\item Accede a la aplicación
			\item Ver citas asociadas del día en la home
			\item Ver todas las citas en el calendario
		\end{enumerate}\\
		\textbf{Postcondición}        & Se muestran las citas \\
		\textbf{Excepciones}          & Error al obtener los datos \\
		\textbf{Importancia}          & Alta \\
		\bottomrule
	\end{tabularx}
	\caption{CU-7 Gestión citas.}
\end{table}

\begin{table}[p]
	\centering
	\begin{tabularx}{\linewidth}{ p{0.21\columnwidth} p{0.71\columnwidth} }
		\toprule
		\textbf{CU-8}    & \textbf{Crear cita}\\
		\toprule
		\textbf{Versión}              & 1.0    \\
		\textbf{Autor}                & {\nombre} \\
		\textbf{Requisitos asociados} & RF-1.8, RF-4.1 \\
		\textbf{Descripción}          & Permite crear una cita\\
		\textbf{Precondición}         & Firebase está disponible \\
		\textbf{Acciones}             &
		\begin{enumerate}
			\def\labelenumi{\arabic{enumi}.}
			\tightlist
			\item Accede a la aplicación
			\item Ir al calendario
			\item Pulsar botón (+) para crear cita
			\item Introducir los datos
			\item Aceptar
		\end{enumerate}\\
		\textbf{Postcondición}        & Se crea la cita \\
		\textbf{Excepciones}          & Error al obtener los datos \\
		\textbf{Importancia}          & Alta \\
		\bottomrule
	\end{tabularx}
	\caption{CU-8 Crear cita.}
\end{table}


\begin{table}[p]
	\centering
	\begin{tabularx}{\linewidth}{ p{0.21\columnwidth} p{0.71\columnwidth} }
		\toprule
		\textbf{CU-9}    & \textbf{Editar cita}\\
		\toprule
		\textbf{Versión}              & 1.0    \\
		\textbf{Autor}                & {\nombre} \\
		\textbf{Requisitos asociados} & RF-1.9 RF-4.4 \\
		\textbf{Descripción}          & Permite editar una cita\\
		\textbf{Precondición}         & Firebase está disponible \\
		\textbf{Acciones}             &
		\begin{enumerate}
			\def\labelenumi{\arabic{enumi}.}
			\tightlist
			\item Accede a la aplicación
			\item Ir al calendario
			\item Pulsar sobre una cita
			\item Cambiar los datos
			\item Aceptar
		\end{enumerate}\\
		\textbf{Postcondición}        & Se modifica la cita \\
		\textbf{Excepciones}          & Error al obtener los datos \\
		\textbf{Importancia}          & Alta \\
		\bottomrule
	\end{tabularx}
	\caption{CU-9 Editar cita.}
\end{table}

\begin{table}[p]
	\centering
	\begin{tabularx}{\linewidth}{ p{0.21\columnwidth} p{0.71\columnwidth} }
		\toprule
		\textbf{CU-10}    & \textbf{Cancelar cita}\\
		\toprule
		\textbf{Versión}              & 1.0    \\
		\textbf{Autor}                & {\nombre} \\
		\textbf{Requisitos asociados} & RF-1.10, RF-2.4, RF-3.4, RF-4.2 \\
		\textbf{Descripción}          & Permite cancelar una cita\\
		\textbf{Precondición}         & Firebase está disponible \\
		\textbf{Acciones}             &
		\begin{enumerate}
			\def\labelenumi{\arabic{enumi}.}
			\tightlist
			\item Accede a la aplicación
			\item Ir al calendario
			\item Pulsar sobre una cita
			\item Pulsar sobre cancelar cita
		\end{enumerate}\\
		\textbf{Postcondición}        & Se cancela la cita \\
		\textbf{Excepciones}          & Error al obtener los datos \\
		\textbf{Importancia}          & Alta \\
		\bottomrule
	\end{tabularx}
	\caption{CU-10 Cancelar cita.}
\end{table}


\begin{table}[p]
	\centering
	\begin{tabularx}{\linewidth}{ p{0.21\columnwidth} p{0.71\columnwidth} }
		\toprule
		\textbf{CU-11}    & \textbf{Listar hospitales}\\
		\toprule
		\textbf{Versión}              & 1.0    \\
		\textbf{Autor}                & {\nombre} \\
		\textbf{Requisitos asociados} & RF-5.9 \\
		\textbf{Descripción}          & Permite ver los hospitales \\
		\textbf{Precondición}         & Firebase está disponible \\
		\textbf{Acciones}             &
		\begin{enumerate}
			\def\labelenumi{\arabic{enumi}.}
			\tightlist
			\item Accede a la aplicación como admin
			\item Pulsar sobre hospitales
		\end{enumerate}\\
		\textbf{Postcondición}        & Se obtienen los hospitales\\
		\textbf{Excepciones}          & Error al obtener los datos \\
		\textbf{Importancia}          & Alta \\
		\bottomrule
	\end{tabularx}
	\caption{CU-11 Listar hospitales.}
\end{table}

\begin{table}[p]
	\centering
	\begin{tabularx}{\linewidth}{ p{0.21\columnwidth} p{0.71\columnwidth} }
		\toprule
		\textbf{CU-12}    & \textbf{Listar tratamientos}\\
		\toprule
		\textbf{Versión}              & 1.0    \\
		\textbf{Autor}                & {\nombre} \\
		\textbf{Requisitos asociados} & RF-5.13 \\
		\textbf{Descripción}          & Permite ver los tratamientos \\
		\textbf{Precondición}         & Firebase está disponible \\
		\textbf{Acciones}             &
		\begin{enumerate}
			\def\labelenumi{\arabic{enumi}.}
			\tightlist
			\item Accede a la aplicación como admin
			\item Pulsar sobre tratamientos
		\end{enumerate}\\
		\textbf{Postcondición}        & Se obtienen los tratamientos\\
		\textbf{Excepciones}          & Error al obtener los datos \\
		\textbf{Importancia}          & Alta \\
		\bottomrule
	\end{tabularx}
	\caption{CU-12 Listar tratamientos.}
\end{table}


\begin{table}[p]
	\centering
	\begin{tabularx}{\linewidth}{ p{0.21\columnwidth} p{0.71\columnwidth} }
		\toprule
		\textbf{CU-13}    & \textbf{Crear hospital}\\
		\toprule
		\textbf{Versión}              & 1.0    \\
		\textbf{Autor}                & {\nombre} \\
		\textbf{Requisitos asociados} & RF-5.10 \\
		\textbf{Descripción}          & Permite crear un hospital \\
		\textbf{Precondición}         & Firebase está disponible \\
		\textbf{Acciones}             &
		\begin{enumerate}
			\def\labelenumi{\arabic{enumi}.}
			\tightlist
			\item Accede a la aplicación como admin
			\item Pulsar sobre hospitales
			\item Pulsar sobre el botón (+)
			\item Introducir los datos
			\item Aceptar
		\end{enumerate}\\
		\textbf{Postcondición}        & Se guarda el hospital\\
		\textbf{Excepciones}          & Error al obtener los datos \\
		\textbf{Importancia}          & Alta \\
		\bottomrule
	\end{tabularx}
	\caption{CU-13 Crear hospital.}
\end{table}

\begin{table}[p]
	\centering
	\begin{tabularx}{\linewidth}{ p{0.21\columnwidth} p{0.71\columnwidth} }
		\toprule
		\textbf{CU-14}    & \textbf{Crear tratamiento}\\
		\toprule
		\textbf{Versión}              & 1.0    \\
		\textbf{Autor}                & {\nombre} \\
		\textbf{Requisitos asociados} & RF-5.14 \\
		\textbf{Descripción}          & Permite crear un tratamiento \\
		\textbf{Precondición}         & Firebase está disponible \\
		\textbf{Acciones}             &
		\begin{enumerate}
			\def\labelenumi{\arabic{enumi}.}
			\tightlist
			\item Accede a la aplicación como admin
			\item Pulsar sobre tratamientos
			\item Pulsar sobre el botón (+)
			\item Introducir los datos
			\item Aceptar
		\end{enumerate}\\
		\textbf{Postcondición}        & Se guarda el tratamiento\\
		\textbf{Excepciones}          & Error al obtener los datos \\
		\textbf{Importancia}          & Alta \\
		\bottomrule
	\end{tabularx}
	\caption{CU-14 Crear tratamiento.}
\end{table}

\begin{table}[p]
	\centering
	\begin{tabularx}{\linewidth}{ p{0.21\columnwidth} p{0.71\columnwidth} }
		\toprule
		\textbf{CU-15}    & \textbf{Editar hospital}\\
		\toprule
		\textbf{Versión}              & 1.0    \\
		\textbf{Autor}                & {\nombre} \\
		\textbf{Requisitos asociados} & RF-5.11 \\
		\textbf{Descripción}          & Permite crear un hospital \\
		\textbf{Precondición}         & Firebase está disponible \\
		\textbf{Acciones}             &
		\begin{enumerate}
			\def\labelenumi{\arabic{enumi}.}
			\tightlist
			\item Accede a la aplicación como admin
			\item Pulsar sobre hospitales
			\item Pulsar sobre un hospital
			\item Introducir los datos
			\item Aceptar
		\end{enumerate}\\
		\textbf{Postcondición}        & Se guarda el hospital\\
		\textbf{Excepciones}          & Error al obtener los datos \\
		\textbf{Importancia}          & Alta \\
		\bottomrule
	\end{tabularx}
	\caption{CU-15 Editar hospital.}
\end{table}

\begin{table}[p]
	\centering
	\begin{tabularx}{\linewidth}{ p{0.21\columnwidth} p{0.71\columnwidth} }
		\toprule
		\textbf{CU-16}    & \textbf{Editar tratamiento}\\
		\toprule
		\textbf{Versión}              & 1.0    \\
		\textbf{Autor}                & {\nombre} \\
		\textbf{Requisitos asociados} & RF-5.15 \\
		\textbf{Descripción}          & Permite editar un tratamiento \\
		\textbf{Precondición}         & Firebase está disponible \\
		\textbf{Acciones}             &
		\begin{enumerate}
			\def\labelenumi{\arabic{enumi}.}
			\tightlist
			\item Accede a la aplicación como admin
			\item Pulsar sobre tratamientos
			\item Pulsar sobre un tratamiento
			\item Introducir los datos
			\item Aceptar
		\end{enumerate}\\
		\textbf{Postcondición}        & Se guarda el tratamiento\\
		\textbf{Excepciones}          & Error al obtener los datos \\
		\textbf{Importancia}          & Alta \\
		\bottomrule
	\end{tabularx}
	\caption{CU-16 Editar tratamiento.}
\end{table}

\begin{table}[p]
	\centering
	\begin{tabularx}{\linewidth}{ p{0.21\columnwidth} p{0.71\columnwidth} }
		\toprule
		\textbf{CU-17}    & \textbf{Eliminar hospital}\\
		\toprule
		\textbf{Versión}              & 1.0    \\
		\textbf{Autor}                & {\nombre} \\
		\textbf{Requisitos asociados} & RF-5.12 \\
		\textbf{Descripción}          & Permite eliminar un hospital \\
		\textbf{Precondición}         & Firebase está disponible \\
		\textbf{Acciones}             &
		\begin{enumerate}
			\def\labelenumi{\arabic{enumi}.}
			\tightlist
			\item Accede a la aplicación como admin
			\item Pulsar sobre hospitales
			\item Pulsar sobre un hospital
			\item pulsar sobre la papelera
			\item Aceptar
		\end{enumerate}\\
		\textbf{Postcondición}        & Se elimina el hospital\\
		\textbf{Excepciones}          & Error al obtener los datos \\
		\textbf{Importancia}          & Alta \\
		\bottomrule
	\end{tabularx}
	\caption{CU-17 Eliminar hospital.}
\end{table}

\begin{table}[p]
	\centering
	\begin{tabularx}{\linewidth}{ p{0.21\columnwidth} p{0.71\columnwidth} }
		\toprule
		\textbf{CU-18}    & \textbf{Eliminar tratamiento}\\
		\toprule
		\textbf{Versión}              & 1.0    \\
		\textbf{Autor}                & {\nombre} \\
		\textbf{Requisitos asociados} & RF-5.17 \\
		\textbf{Descripción}          & Permite eliminar un tratamiento \\
		\textbf{Precondición}         & Firebase está disponible \\
		\textbf{Acciones}             &
		\begin{enumerate}
			\def\labelenumi{\arabic{enumi}.}
			\tightlist
			\item Accede a la aplicación como admin
			\item Pulsar sobre tratamientos
			\item Pulsar sobre la papelera
			\item Aceptar
		\end{enumerate}\\
		\textbf{Postcondición}        & Se elimina el tratamiento\\
		\textbf{Excepciones}          & Error al obtener los datos \\
		\textbf{Importancia}          & Alta \\
		\bottomrule
	\end{tabularx}
	\caption{CU-18 Eliminar tratamiento.}
\end{table}
\apendice{Especificación de diseño}

\section{Introducción}
En este anexo vamos a definir en base a los objetivos y especificaciones que datos va a manejar la aplicación, que estructura, arquitectura y diseño tendrá.

\section{Diseño de datos}
Nuestra aplicación cuenta con un diseño de base de datos \textbf{NoSQL}, lo que nos permite mediante la flexibilidad de la estructura de datos adaptarnos rápidamente a los requisitos y necesidades del negocio. Mediante los documentos que darán vida a nuestras entidades estructurados con formato JSON, nos permitirán crear una aplicación donde el usuario no esté relacionado directamente como sucede en las bases de datos relacionales y nos permita mayor rendimiento cuando la aplicación vaya creciendo.

A continuación, vamos a definir el diseño de datos con el que cuenta nuestra aplicación. Para poder entenderlo más rápidamente, vamos a representar cada elemento de nuestra aplicación como si fuera una una entidad relacional.
\imagenTam{diseño-datos}{Diseño de datos actores}{.7}
 
Otras entidades que necesitaremos para nuestra aplicación son:
\imagen{entidades}{Entidades}

\section{Diseño procedimental}
En este punto vamos a definir los aspectos más relevantes de la aplicación.

\subsection{Envío invitación}
Planteamos la idea de que en esta aplicación los usuarios que se puedan registrar en la aplicación sea a partir de invitaciones que llegarán al correo electrónico. Para poder recibir una invitación se tiene que tener en cuenta lo siguiente:
\begin{itemize}
	\item Administrador puede invitar a nuevos administradores o gestores de casos.
	\item Gestor de casos puede invitar a nuevos gestores de casos o pacientes.
	\item Paciente puede invitar a que se registre su cuidador personal.
\end{itemize}

\imagen{diagrama-secuencia-invitación}{Diagrama de secuencia invitar Usuario}

\subsection{Registro usuario invitado}
Una vez el usuario ha recibido la invitación, este se registrará mediante OTP, que es una contraseña de un solo uso y posteriormente rellenará los datos del registro. Una vez guardado, accederá a la pantalla principal.

\imagenTam{diagrama-secuencia-registro}{Diagrama de secuencia registro usuario}{.8}

\subsection{Crear cita gestor de casos}
lL gestión de las citas por parte del gestor de casos. En ella mostramos el diagrama de secuencia con la creación y borrado de una cita donde se notificará a los pacientes.
\imagenTam{diagrama-secuencia-gestor-casos}{Gestión citas Gestor de casos}{.8}

\section{Diseño arquitectónico}
Al elegir hacer la aplicación con Flutter nos vimos ante la necesidad de realizar una arquitectura limpia y por esa misma razón escogimos usar MVVM (Model-View-ViewModel) estructurando nuestras carpetas de la siguiente forma:
\begin{itemize}
	\item nuestro model contendrá las entidades que controla la aplicación
	\item nuestro view contendrá las vistas de la aplicación
	\item nuestro viewmodel hará de intermediario entre la parte de los datos y la interfaz.
\end{itemize}

%\imagenTam{mvvm-pattern}{MVVM}{0.8}



\apendice{Documentación técnica de programación}

\section{Introducción}

Este anexo va orientado a describir y documentar la información necesaria sobre la poder realizar la instalación de las herramientas necesarias para el desarrollo, que estructura tiene la aplicación, como compilarlo, que servicios se necesitan.

- Cuenta firebase
- Generar dotenv
- Descargar el proyecto
- Comandos para arrancar el proyecto

\section{Estructura de directorios}
La estructura de nuestro proyecto está orientado en la arquitectura MVVM y será la siguiente;
\imagen{mvvm-clean-arquitecture}{mvvm}

\begin{itemize}
	\item \textcolor{darkblue}{\url{/.github}}: carpeta contiene los actions para los despliegues CI/CD en web y android al realizar push sobre la rama \textbf{main}.
	\item \textcolor{darkblue}{\url{/android}}: carpeta que contiene el código para la aplicación android.
	\item \textcolor{darkblue}{\url{/assets}}: carpeta que contendrá las imágenas y  fichero dotenv para las variables de entorno.
	\item \textcolor{darkblue}{\url{/doc}}: carpeta que contendrá la memoria y anexos.
	\item \textcolor{darkblue}{\url{/lib}}: directorio donde se encuentra el código fuente del proyecto.
		\begin{itemize}
			\item \textcolor{darkblue}{\url{main.dart}}: fichero principal del proyecto.
			\item \textcolor{darkblue}{\url{/app}}: contendrá la estructura de carpetas de nuestro Model-View-ViewModel
			\begin{itemize}
				\item \textcolor{darkblue}{\url{/data}}: contiene las implementaciones a los repositorios
				\item \textcolor{darkblue}{\url{/domain}}: contiene los modelos y las interfaces de los repositorios
				\item \textcolor{darkblue}{\url{/presententation}}: contiene los modulos que representan las vistas de nuestra aplicación
			\end{itemize}
			\item \textcolor{darkblue}{\url{/l10n}}: contrendrá los archivos para la internacionalización de la aplicación
		\end{itemize}
	\item \textcolor{darkblue}{\url{/web}}: carpeta que contiene el código para la aplicación web
	\item \textcolor{darkblue}{\url{analysis_options.yaml}}
	\item \textcolor{darkblue}{\url{build.yaml}}
	\item \textcolor{darkblue}{\url{pubspec.yaml}}: archivo que contendrá los paquetes necesarios para el correcto funcionamiento de la aplicación
	\item \textcolor{darkblue}{\url{l10n.yaml}}: archivo para la internacionalización de la aplicación
\end{itemize}


\section{Manual del programador}
Este manual está pensado para que futuros programadores puedan poner en marcha el proyecto y saber que pasos deben seguir para el correcto funcionamiento.

Cómo primer requisito es instalar Flutter~\cite{flutter}, para ello tendremos que ir a la página oficial de flutter \url{https://docs.flutter.dev/get-started/install} y elegir el sistema donde vamos a realizar la instalación y para que tipo de aplicación queremos desarrollar (en nuestro caso android).
\imagen{flutter-so}{Plataforma desarrollo}

Una vez realizada las instalaciones y configuraciones correspondientes, debemos comprobar que todo ha ido correctamente, para ello utilizaremos la siguiente instrucción en el terminal de \textbf{VS Code} \textbf{flutter doctor}
\imagen{flutter-doctor-command}{Comprobación correcta instalación Flutter}

\subsection{Visual Studio Code}
Para ello debemos de tener instalado VS Code, en caso contrario descargar de la página oficial \url{https://code.visualstudio.com/}.

Una vez instalado añadimos la extensión de Flutter
\imagen{flutter-extension}{Añadir Flutter a Visual Studio Code}

\subsection{Flutter SDK}
Siguiendo los pasos que indican en la guía de Flutter, nosotros escogimos instalar Flutter SDK desde VS Code, para ello debemos hacer:
\begin{enumerate}
	\item Abrir \textbf{VS Code}
	\item Pulsar \textbf{Ctrl}+\textbf{Shift}+\textbf{P} para abrir la paleta de comandos
	\item Seleccionar \textbf{New Project}
		\begin{itemize}
			\item Si no lo tienes descargado, pulsa en \textbf{Download SDK}.
			\item Si lo tienes descargado, pulsa \textbf{Locate SDK}
		\end{itemize}
	\item Cuando te pregunte que plantilla usar, hay que ignorarlo.
\end{enumerate}

\imagen{flutter-command-palette}{Paleta de comandos VS Code}

\subsection{Android Studio}
Para ello si no lo tenemos instalado, tendremos que realizar la instalación de \textbf{Android Studio} desde \url{https://developer.android.com/studio/install#windows} en mi caso lo instalo para Windows. Una vez instalado deberemos instalar el plugin de Flutter para Intellij \url{https://plugins.jetbrains.com/plugin/9212-flutter}

Para poder crear aplicación con Flutter, necesitamos comprobar que tenemos instalados los siguientes componentes de Android:
\begin{itemize}
	\item Android SDK Platform, API 34.0.5
	\item Android SDK Command-line Tools
	\item Android SDK Build-Tools
	\item Android SDK Platform-Tools
	\item Android Emulator
\end{itemize}

\subsubsection{Configuracion Android Studio}
Realizar la siguiente configuración de android studio.
\begin{enumerate}
	\item Habilite la aceleración de VM en su computadora de desarrollo.
	\item Inicie \textbf{Android Studio}.
	\item Vaya al diálogo de Configuración para ver el \textbf{Administrador de SDK}.
	\begin{itemize}
		\item Si tiene un proyecto abierto, vaya a \textbf{Herramientas} > \textbf{Administrador de dispositivos}.
		\item Si se muestra el diálogo de bienvenida, haga clic en el icono pulsar \textbf{más opciones} que sigue al botón \textbf{abrir} y seleccione \textbf{administrador de dispositivos} desde el menú desplegable.
	\end{itemize}
	\item Haga clic en \textbf{virtual}.
	\item Haga clic en \textbf{crear dispositivo}. El diálogo de configuración de dispositivo virtual se muestra.
	\item Seleccione el tipo de dispositivo bajo la categoría.
	\item Seleccione una definición de dispositivo. Puede explorar o buscar por el dispositivo.
	\item Haga clic en \textbf{siguiente}.
	\item Haga clic en \textbf{imagenes x86}.
	\item Haga clic en una imagen del sistema para la versión de Android que desee emular.
	\begin{itemize}
		\item Si la imagen deseada tiene un icono de descarga al lado del nombre de lanzamiento, haga clic en él.
		 \item Cuando complete la descarga, haga clic en \textbf{finalizar}.
	\end{itemize}
	\item Haga clic en \textbf{siguiente}. La configuración del dispositivo virtual se muestra en su paso de verificación.
	\item Para renombrar el \textbf{Dispositivo Virtual de Android (AVD)}, cambie el valor en el cuadro de texto \textbf{nombre AVD}.
	\item Haga clic en \textbf{mostrar configuraciones avanzadas} y desplace hasta \textbf{rendimiento emulado}.
	\item Desde el menú desplegable \textbf{gráficos}, seleccione \textbf{Hardware - GLES 2.0}. (Esto habilita la aceleración de hardware y mejora el rendimiento de renderizado).
	\item Verifique la configuración de su AVD. Si todo está correcto, haga clic en \textbf{finalizar}.
	\item En el diálogo del Administrador de dispositivos, haga clic en el icono \textbf{ejecutar} al lado derecho del AVD deseado. El emulador se inicia y muestra la pantalla predeterminada para la versión de Android OS y dispositivo seleccionados.
\end{enumerate}

\subsubsection{Fichero variables de entorno}
Antes de generar el fichero, debemos seguir los pasos descritos en el paquete de \textbf{Flutter} \textbf{encrypt} \url{https://pub.dev/packages/encrypt}:
\begin{itemize}
	\item Activar el paquete \fboxrule=1pt\fboxsep=4pt\fcolorbox{darkblue}{cgoLight}{pub global activate encrypt}
\end{itemize}

Para generar las claves seguras de manera aleatoria utilizaremos la siguiente instrucción \fboxrule=1pt\fboxsep=4pt\fcolorbox{darkblue}{cgoLight}{secure-random}. Con ella podremos generar claves que nos servirán para encriptar de manera segura los datos tratados por nuestra aplicación.
\imagen{flutter-encrypt-secure-random}{Encrypt secure random}

Este fichero lo tenemos que generar dentro de la carpeta \textbf{assets} con el nombre \textbf{dotenv}, ya que está ignorado en las subidas a \textbf{Github}. El fichero contendrá unas claves necesarias para el correcto funcionamiento de la aplicación:
\begin{itemize}
	\item \textbf{ENCRYPT\_KEY} esta clave la generaremos con la instrucción anterior y cuya longitud debe de ser de 32
	\item \textbf{ENCRYPT\_IV }esta clave la generaremos con la instrucción anterior y cuya longitud debe de ser de 12
	\item \textbf{USER\_EMAIL}=redelaubutfg@gmail.com
	\item \textbf{USER\_PASS}=qdmucyqwausbluzw
	\item \textbf{VAPID\_KEY} esta clave la necesitamos para poder obtener notificaciones push en nuestra web. Explicado en \textbf{\ref{subsec:firebase} Firebase}
\end{itemize}

\subsubsection{\label{subsec:firebase}Firebase}
Crearemos una cuenta Firebase y agregaremos un proyecto con el nombre \textbf{RedELA}
\begin{enumerate}
\item Agregaremos el proyecto 
\imagen{firebase-agregar-proyecto}{Agregar proyecto}
\item Introduciremos el nombre del proyecto. Pulsaremos \textbf{continuar}
\imagen{firebase-nombre-proyecto}{Agregar proyecto}
\item Deshabilitaremos Google Analytics y pulsaremos \textbf{crear proyecto}
\end{enumerate}

Tras unos segundos tendremos creado nuestro proyecto, ahora necesitamos enlazar ese proyecto con nuestro código descargado desde \textbf{\ref{subsec:github} Github} \url{https://github.com/jmc1005/red-ela}

Una vez creado nuestro proyecto en Firebase, debemos obtener el \textbf{VAPID\_KEY} en \textbf{descripción general} \textbf{configuración del proyecto} vamos a la sección de \textbf{cloud messaging} y al final en \textbf{configuración web} en \textbf{certificacdos push web} encontraremos en \textbf{par de claves} nuestro \textbf{VAPID\_KEY}.

\subsubsection{\label{subsec:github}Github}
Para empezar, necesitamos crear nuestros secrets en Github. Para acceder a \textbf{secrets} debemos a acceder a \textbf{Settings} y luego sobre el menú lateral izquierdo ir a \textbf{Secrets and variables} y pulsaremos sobre actions
\imagen{github-secrets}{Secrets and variables}

Posteriormente, sobre \textbf{Repository secrets} añadiremos nuestras variables necesarias para los despliegues mediante los workflow.
\imagen{github-repository-secrets}{Respository secrets}

En este punto vamos a explicar un poco a cerca de los workflows definidos en nuestro proyecto y además crear los secrets que serán necesarios para nuestro proyecto.

Los workflow son acciones que nos van a permitir desplegar en el hosting de firebase nuestra web y generar nuestras apks 

\subsubsection{Google Play}
De momento sólo estamos haciendo despliegues en plan Prueba Interna, para poder desplegar la aplicación para prueba interna necesitamos generar:
\begin{enumerate}
	\item Modificar nuestro fichero \textbf{pubspec\.yaml} cambiando la versión por uno más es decir, si tenemos \textbf{version: 1.0.3+3} poner \textbf{version: 1.0.4+4}
	\item un \textbf{bundle} con el comando \fboxrule=1pt\fboxsep=4pt\fcolorbox{darkblue}{cgoLight}{flutter build appbundle}, una vez generado el \textbf{bundle} accedemos a nuestra \textbf{consola de google play} \url{https://play.google.com/console/u/0/developers}
	\item pulsar en \textbf{crear nueva versión}
\imagen{google-play-nueva-version} {Nueva versión}
	\item Subir el bundle generado y actualizar la versión, por ejemplo a 4 (1.0.4)
\imagen{google-play-subir-bundle} {Subir bundle}
\imagen{google-play-detalle-version} {Detalle versión}
\end{enumerate}

\section{Compilación}
Para poder compilar el proyecto debemos ejecutar estos comandos \fboxrule=1pt\fboxsep=4pt\fcolorbox{darkblue}{cgoLight}{dart run build\_runner build} y \fboxrule=1pt\fboxsep=4pt\fcolorbox{darkblue}{cgoLight}{flutter pub get}. El primer comando nos generará el código automático de las clases que están con la anotación @freezed en nuestro proyecto. Estas anotaciones nos permiten generar código automático con pocas líneas de código.

\section{Ejecución del proyecto}
La aplicación se puede ejecutar tanto en la web como en un dispositivo móvil. Para ejecutar las pruebas podemos realizarlas tanto en un dispositivo móvil con Android que tenga API 34 o bien accediendo a la web.

También se pueden hacer pruebas en local desde VSCode \fboxrule=1pt\fboxsep=4pt\fcolorbox{darkblue}{cgoLight}{flutter run -d chrome} si queremos que se ejecute en chrome o bien seleccionando un dispositivo como se muestra en la imagen.
\imagen{flutter-device}{Seleccionar dispositivo}

\section{Pruebas del sistema}
Se ha generado una clase que testea mediante \textbf{flutter\_test} Firebase que será el gestione los datos que utiliza la aplicación.

Para poder ejecutar los test deberemos pulsar los iconos que serán un play
\imagen{flutter-test}{Test}

Cuando los tests pasan se verá de la siguiente manera.
\imagen{flutter-test-ok}{Test OK}

\apendice{Documentación de usuario}

\section{Introducción}

\section{Requisitos de usuarios}

\section{Instalación}

\section{Manual del usuario}



\apendice{Anexo de sostenibilización curricular}

\section{Introducción}
El objetivo del proyecto es crear una aplicación que facilite el día a día de los pacientes diagnosticados con ELA (Esclerósis Lateral Amiotrófica)~\cite{wiki:ela}. Esta aplicación está creada con el framework Flutter~\cite{flutter} y con la plataforma Firebase~\cite{firebase} que nos permite gestionar incialmente las citas y tratamientos que tiene le paciente. 

Además, con el desarrollo de la aplicación, se pretende concienciar y dar sostenibilidad al cuidado de personas con ELA haciendo que se comprometan aquellos profesionales de la salud y cuidadores principales ante los cuidados necesarios para hacer la vida más fácil a personas que hayan sido diagnosticadas con la enfermedad.

Para realizar el anexo de sostenibilidad se utiliza como referencia el documento de la CRUE~\cite{crue}

\subsection{Objetivos}
Algunos de los objetivos que podemos destacar son:
\begin{itemize}
\tightlist 
\item Concienciar la importancia sobre la sostenibilidad sobre los cuidados que necesitan los paciente con ELA.
\item Concienciar a crear soluciones sostenibles para el cuidado de pacientes con ELA.
\item Realizar desarrollos que promuevan la responsabilidad para dar solución a pacientes con ELA.
\end{itemize}

\subsection{Estrategias}
Algunas estrategias que se deben de usar son:
\begin{itemize}
\tightlist 
\item Realizar formaciones que fomenten la colaboración entre los profesionales de la salud, pacientes y cuidadores principales.
\item Sensibilizar y dar visibilidad a los pacientes que padezcan esta enfermedad.
\item Realizar desarrollos que promuevan la responsabilidad para dar solución a pacientes con ELA.
\end{itemize}

\subsection{Conclusión}
Sacamos como conclusión que la realización del proyecto aporta a la soluciones a la sociedad facilitando la vida a personas enfermas de ELA, a los cuidadores principales e incluso a los profesionales sanitarios. Además, a nivel medio ambiental reduce el uso de emisiones contaminantes ahorrando en el gasto de papel. 
Con ello se tendrá en cuenta la aceptación, demanda y aceptación por parte del usuario de la aplicación para realizar ajustes y nuevas funcionalidades a lo largo del tiempo


\bibliographystyle{plain}
\bibliography{bibliografiaAnexos}

\end{document}
